% Generated draft: requires mathematical and source review.
\documentclass[11pt]{article}
\usepackage[margin=25mm]{geometry}
\usepackage{fontspec}
\setmainfont{Latin Modern Roman}
\usepackage{amsmath,amssymb,mathtools,unicode-math}
\setmathfont{Latin Modern Math}
\usepackage{enumitem,needspace,xurl,hyperref}
\setlength{\parindent}{0pt}
\setlength{\parskip}{5pt}
\setlength{\emergencystretch}{3em}
\newcommand{\N}{\mathbb{N}}
\newcommand{\Z}{\mathbb{Z}}
\newcommand{\Q}{\mathbb{Q}}
\newcommand{\R}{\mathbb{R}}
\newcommand{\C}{\mathbb{C}}
\newcommand{\F}{\mathbb{F}}
\newcommand{\A}{\mathbb{A}}
\newcommand{\PP}{\mathbb{P}}
\newcommand{\E}{\mathbb{E}}
\newcommand{\eps}{\varepsilon}
\newcommand{\dd}{\,\mathrm{d}}
\begin{document}
\section*{UnsolvedMath ID: 700884}
\textbf{Source problem code:} catalogue entry
\subsection{Definitions and mathematical statement}
Let \(\Pi=(P,\mathcal L,I)\) be a finite affine plane: any two distinct points lie on exactly one line; for each line \(\ell\) and point \(p\) not on \(\ell\), exactly one line through \(p\) is disjoint from \(\ell\); and there exist three noncollinear points. Lines are parallel if they coincide or are disjoint. A collineation is an incidence-preserving permutation of the points and lines. Write \(G=\operatorname{Aut}(\Pi)\) for the full group of affine collineations.

A translation is the identity collineation or a collineation having no fixed affine point and mapping every line to a parallel line. The plane \(\Pi\) is a translation plane if its translations act transitively on \(P\); this action is then sharply transitive.

The action of \(G\) on \(P\) is doubly transitive if
\[
\forall x_1,x_2,y_1,y_2\in P,\qquad
\bigl(x_1\ne x_2\ \text{and}\ y_1\ne y_2\bigr)
\Longrightarrow
\exists g\in G:\quad g(x_1)=y_1,\quad g(x_2)=y_2.
\]
Describe, up to isomorphism of affine incidence structures, all finite translation planes satisfying this condition.

Here ``collineation groups'' is interpreted through the full affine collineation group. For this property, requiring the existence of a doubly transitive subgroup is equivalent. The supplied source labels this classification problem open; that label is supplied curation, not an independently verified assessment.
\subsection{Short English statement}
Classify the finite translation planes whose affine collineations can send any ordered pair of distinct affine points to any other such pair. The problem is attributed to N.\ D.\ Podufalov (1992).
\subsection{Sources}
\begin{itemize}
\item Local UnsolvedMath registry, permanent ID \(700884\), exact source code \texttt{catalogue entry}, title \emph{Kourovka Notebook Problem 12.66}. The supplied record labels the problem \texttt{open}.
\item \emph{The Kourovka Notebook (Unsolved Problems in Group Theory)}, original problem 12.66, 12th issue (1992), proposed by N.\ D.\ Podufalov. The supplied background also identifies the 21st issue edition (2026), arXiv:1401.0300v48. The supplied collection description names E.\ I.\ Khukhro and V.\ D.\ Mazurov as editors.
\item These references and the status assessment are reproduced from the supplied record and have not been independently verified. No separate literature-verification date or literature-source links were supplied.
\end{itemize}
\end{document}
